\documentclass[../../../include/open-logic-section]{subfiles}

\begin{document}

\olfileid{sth}{replacement}{absinf}
\olsection{ప్రతిస్థాపన, ``సంపూర్ణ అనంతత్వం''}

ఇప్పుడు \limofsize{}ను పక్కనపెట్టి, సమితులు దశల్లో ఏర్పడతాయనే ఆలోచనను గంభీరంగా తీసుకునే, ప్రతిస్థాపనకు (అంతర్గత) సమర్థన ఇచ్చే మరో ప్రయత్నాల కుటుంబాన్ని పరిశీలిద్దాం.

పునరావర్తన ప్రక్రియను మొదట వివరించినప్పుడు, ప్రతి దశలో ఏమి జరుగుతుందో చెప్పే కొన్ని సూత్రాలను ఇచ్చాం. అవి \stageshier, \stagesord, \stagesacc. తరువాత దశల సంఖ్య గురించి కొన్ని సూత్రాలను చేర్చాం: సమితి-నిర్మాణ ప్రక్రియ ఎన్నడూ ముగియదని \stagessucc{} చెప్పింది; ప్రక్రియ అనంతవ దశ వరకు సాగుతుందని \stagesinf{} చెప్పింది.

ఈ చివరి రెండు సూత్రాలు మరింత విస్తృతమైన ఒక సూత్రం నుండి వస్తాయని సూచించడం సమంజసమే:
\begin{enumerate}
	\item[] \stagesinex. దశలు సంపూర్ణంగా అనంత సంఖ్యలో ఉన్నాయి; స్థాయిక్రమం సాధ్యమైనంత ఎత్తుగా ఉంటుంది.
\end{enumerate}
స్పష్టంగా ఇది అనౌపచారిక సూత్రం. సంచిత-పునరావర్తన సమితి భావన నుండి ఇది వెంటనే \emph{అనుసరించకపోయినా}, దానితో \emph{పొసగుతుంది} అనిపిస్తుంది. కనీసం \limofsize{}కు భిన్నంగా, సమితులు దశలవారీగా ఏర్పడతాయనే ఆలోచనను ఇది నిలుపుతుంది.

ఇప్పుడు \stagesinex{} ఆధారంగా ప్రతిస్థాపనను సమర్థించాలనే ఆశ. ఇది ఎలా సాధ్యమో చూద్దాం.

\olref[ordinals][idea]{sec}లో, (భావనాపరంగా) $\omega+\omega$కు సమరూపంగా ఉండవలసిన ఒక సుక్రమాన్ని సులభంగా నిర్మించవచ్చని చూశాం. మరో మాటలో, (భావనాపరంగా) $\omega+\omega$ క్రమరకం గల దశలవారీ పునరావర్తన ప్రక్రియను సులభంగా ఊహించవచ్చు. కాబట్టి \stagesinex{}ను అంగీకరిస్తే, స్థాయిక్రమంలో కనీసం $\omega+\omega$వ దశ, అంటే $V_{\omega+\omega}$, ఉందని తప్పకుండా అంగీకరించాలి; స్థాయిక్రమం అంతదూరం \emph{సాగగలదు} కదా.

ఈ ఆలోచనను ఇలా సాధారణీకరించవచ్చు: ఏ సుక్రమానికైనా, పునరావర్తన స్థాయిక్రమ నిర్మాణ ప్రక్రియ కనీసం ఆ సుక్రమం అంతదూరం సాగాలి. \olref[ordinals][ordtype]{thmOrdinalRepresentation}ను \emph{స్వయంసిద్ధ సూత్రంగా} తీసుకుంటే దీనికి హామీ ఇవ్వవచ్చు. అప్పుడు ఏ సుక్రమమైనా ఒక వాన్ నాయ్‌మన్ క్రమసంఖ్యకు సమరూపమని తెలుస్తుంది. ప్రతి వాన్ నాయ్‌మన్ క్రమసంఖ్య తన స్థాయికే సమానమవుతుంది కనుక, మన సమితి సిద్ధాంతంలో ఒక సుక్రమాన్ని వివరించగలిగినప్పుడల్లా సమితి-నిర్మాణ పునరావర్తన ప్రక్రియ ఆ సుక్రమాన్ని దాటి సాగాలని \olref[ordinals][ordtype]{thmOrdinalRepresentation} చెబుతుంది.

ఈ ఆలోచన \stagesinex{} యొక్క పర్యవసానంలాగే కనిపిస్తుంది. దురదృష్టవశాత్తు, దీని నుండి ప్రతిస్థాపనను పొందాలంటే ఒక సరళమైన సాంకేతిక అడ్డంకి ఉంది: ప్రతిస్థాపన, \olref[ordinals][ordtype]{thmOrdinalRepresentation}కంటే ఖచ్చితంగా బలమైనది. (ఈ విషయాన్ని \citet[\S13.2]{Potter2004} గమనించారు; \olref[sth][replacement][finiteaxiomatizability]{sec}లో మనం నిరూపిస్తాం.)

కాబట్టి \stagesinex{} నుండి ప్రతిస్థాపన రావాలంటే, స్థాయిక్రమం ఏ సుక్రమాన్నైనా దాటి సాగుతుందని అది \emph{మాత్రమే} చెప్పకూడదు. దానికంటే బలమైన వాదన కావాలి. అందుకోసం \cite{Shoenfield:AST} ఈ ఆలోచనకు చాలా సహజమైన బలవర్ధనను ప్రతిపాదించారు: స్థాయిక్రమానికి ఏ సమితీ \emph{సహాంత్యమైనది} కాదు.\footnote{గోడెల్ కూడా ఇలాంటి ఆలోచనను ప్రతిపాదించినట్లు కనిపిస్తుంది; \citet[p.~223]{Potter2004} చూడండి. గోడెల్, \citeauthor{Shoenfield:AST} గురించి చర్చకు \citet[90--5]{Incurvati2020} చూడండి.} మరింత వివరంగా: $\tau$ సమితులను స్థాయిక్రమ దశలకు పంపే చిత్రణ అయితే, ఏ సమితి $A$ యొక్క $\tau$-బింబమూ స్థాయిక్రమాన్ని పూర్తిగా చేరదు. మరో మాటలో (పథకరూపంలో):
\begin{enumerate}
	\item[] \stagescofin. $A$ సమితి కాగా ప్రతి $x \in A$కు $\tau(x)$ ఒక దశ అయితే, ప్రతి $x \in A$కు $\tau(x)$ తరువాత వచ్చే ఒక దశ ఉంటుంది.
\end{enumerate}
$\ZF$లో \stagescofin{}కు తగిన ఔపచారిక రూపం నిరూపించబడుతుందని స్పష్టమే. తిరిగి, \stagescofin{} ప్రతిస్థాపనను సమర్థిస్తుందని అనౌపచారికంగా వాదించవచ్చు.\footnote{\stagescofin{} యొక్క ఏదో ఔపచారిక రూపంతో ప్రతిస్థాపనను నిరూపించడం కష్టమే; ఎందుకంటే $\Z$ ఒక్కటితో దశలను నిర్వచించడానికి తగిన బలం లేదు, కాబట్టి \stagescofin{}ను ఎలా ఔపచారికీకరించాలో స్పష్టం కాదు. అయితే \olref[sth][replacement][extrinsic]{sec}లో చర్చించినట్లుగా $\LT$ యొక్క ఏదో విస్తరణలో పనిచేయడం ఒక మార్గం.} $(\forall x \in A)\lexists![y][\phi(x,y)]$ అనుకుందాం. అప్పుడు ప్రతి $x \in A$కు, $\phi(x,y)$ అయ్యే $y$ను $\sigma(x)$గా, $\sigma(x)$ మొదట ఏర్పడే దశను $\tau(x)$గా ఉంచుదాం. \stagescofin{} ప్రకారం $(\forall x \in A)\tau(x)\in V$ అయ్యే దశ $V$ ఒకటి ఉంది. ఇప్పుడు ప్రతి $\tau(x) \in V$, $\sigma(x) \subseteq \tau(x)$ కాబట్టి, వేరుచేయడం ద్వారా $\Setabs{y
\in V}{(\exists x \in A)\sigma(x) = y} = \Setabs{y}{(\exists x \in
A)\phi(x,y)}$ పొందవచ్చు.

\begin{prob}
	$\ZF$లో \stagescofin{}ను ఔపచారికీకరించండి.
\end{prob}

ఈ విధంగా \stagescofin{} ప్రతిస్థాపనను సమర్థిస్తుంది. అలాగే \stagesinex{} నుండి \stagescofin{} వస్తుందనడం కనీసం సంభావ్యమే. ఎందుకంటే \stagescofin{} విఫలమైతే ఏమవుతుందో అనుకుందాం. అప్పుడు స్థాయిక్రమానికి ఏదో సమితి~$A$ సహాంత్యమైనది; అంటే ఏ దశ~$S$కైనా $S \in \tau(x)$ అయ్యే ఏదో $x \in A$ ఉండేలా $\tau$ చిత్రణ ఉంటుంది. ఆ సందర్భంలో స్థాయిక్రమపు ఊహించిన ``సంపూర్ణ అనంతత్వాన్ని'' ఒక విధంగా పట్టుకోగలం: $A$పై $\tau$ యొక్క విలువల సముదాయం దాన్ని \emph{పూర్తిగా చేరుతుంది}. అది స్థాయిక్రమం ``సంపూర్ణంగా అనంతం'' అనే ఆలోచనను దెబ్బతీస్తుంది. వ్యతిరేక రూపంలో: \stagesinex{} నుండి \stagescofin{} వస్తుంది; అది ప్రతిస్థాపనను సమర్థిస్తుంది.

ఇది ప్రతిస్థాపనకు అంతర్గత సమర్థన ఇచ్చే నిజంగా ఆశాజనకమైన ప్రయత్నం. అయితే చివరికి అది ఫలిస్తుందో లేదో నిర్ణయించాల్సింది మీరే.

\end{document}
